\documentclass[11pt,a4paper]{article}

\usepackage[utf8]{inputenc}
\usepackage[T1]{fontenc}
\usepackage[english]{babel}
\usepackage{amsmath,amssymb,mathtools,bm}
\usepackage{geometry}
\usepackage{booktabs}
\usepackage{array}
\usepackage{microtype}
\usepackage{hyperref}
\geometry{margin=2cm}

\newcommand{\Grad}{\nabla}
\newcommand{\tr}{\operatorname{tr}}
\newcommand{\dev}{\operatorname{dev}}
\newcommand{\Id}{\textbf{I}}
\newcommand{\T}{\mathsf{T}}
\newcommand{\dd}{\text{d}}

\title{Manufactured Dynamic Solution for a Compressible Neo-Hookean Solid}
\author{Pablo Castrillo}
\date{}

\begin{document}
\maketitle
%\tableofcontents

\section{Dynamic problem in the reference configuration}

The solid is formulated in the reference configuration. Material coordinates are written as $\textbf{X}$ and the current position is:
%
\begin{equation}
\textbf{x}(\textbf{X},t)=\textbf{X}+\textbf{u}(\textbf{X},t),
\end{equation}
%
where $\textbf{u}$ is the displacement field.

The balance of linear momentum in material form is:
%
\begin{equation}\label{eq_momentum}
\rho_0\,\ddot{\textbf{u}}=\Grad_0\cdot\textbf{P}+\textbf{f}_0,
\end{equation}
%
where $\rho_0$ is the reference mass density, $\textbf{P}$ is the first Piola--Kirchhoff stress, and $\textbf{f}_0$ is the body force per unit reference volume. Therefore, for a prescribed exact displacement field,
%
\begin{equation}\label{eq_manufacturedForce}
\textbf{f}_0(\textbf{X},t)
=\rho_0\ddot{\textbf{u}}^{\text{ex}}
-\Grad_0\cdot\textbf{P}(\textbf{u}^{\text{ex}}).
\end{equation}

All external boundaries are Dirichlet boundaries,
%
\begin{equation}\label{eq_allDirichlet}
\textbf{u}=\textbf{u}^{\text{ex}}
\qquad\text{on }\partial\Omega.
\end{equation}
%
No traction boundary condition is required for this manufactured problem.

\section{Neo-Hookean model with isochoric--volumetric split}

\subsection{Finite-deformation kinematics}

The deformation gradient \textbf{F}, Jacobian $J$, right Cauchy--Green tensor \textbf{C}, left Cauchy--Green tensor \textbf{B}, and Green--Lagrange strain \textbf{E} are
%
\begin{align}
\textbf{F}&=\dfrac{\partial\,\textbf{x}}{\partial\,\textbf{X}}
=\Id+\Grad_0\textbf{u},
&J&=\det(\textbf{F})>0,
\\
\textbf{C}&=\textbf{F}^{\text{T}}\textbf{F},
&\textbf{B}&=\textbf{F}\textbf{F}^{\text{T}},
\\
\textbf{E}&=\dfrac12(\textbf{C}-\Id).
\end{align}


\subsection{Strain-energy density and derivation of the stress measures}

The strain-energy density is additively split into isochoric and volumetric parts,
%
\begin{equation}
\Psi=\Psi_{\text{iso}}+\Psi_{\text{vol}},
\end{equation}
%
with
%
\begin{equation}
\Psi_{\text{iso}}
=\dfrac{\mu}{2}(\overline I_1-3),
\qquad
\Psi_{\text{vol}}
=\dfrac{\kappa}{4}\left(J^2-1-2\ln J\right),
\end{equation}
%
where\footnote{In terms of Young's modulus $E$ and Poisson's ratio $\nu$,
$$
\mu=\dfrac{E}{2(1+\nu)},
\qquad
\kappa=\dfrac{E}{3(1-2\nu)}=\kappa=\lambda+\dfrac23\mu\longleftarrow \text{Lam\'e parameters}.
$$
}
%
\begin{equation}
\overline I_1=J^{-2/3}I_1\quad\text{and}\quad I_1=\tr\textbf{C}=\tr\textbf{B}.
\end{equation}

Thus
%
\begin{equation}\label{eq_energy}
\Psi(\textbf{F})
=\dfrac{\mu}{2}\left(J^{-2/3}I_1-3\right)
+\dfrac{\kappa}{4}\left(J^2-1-2\ln J\right).
\end{equation}
%

\paragraph{Derivation of the stress measures} Since
%
\begin{eqnarray}
\det\left( \textbf{C}\right)=\det\left( \textbf{F}^{\text{T}}\textbf{F}\right)=\det\left( \textbf{F}^{\text{T}}\right)\det\left(\textbf{F}\right)=J^2\underbrace{\Rightarrow}_{J>0} \boxed{J=\sqrt{\det\left( \textbf{C}\right)}},
\end{eqnarray}
%
\begin{eqnarray}
\boxed{\dfrac{\partial (J^{-2/3})}{\partial J}
	=-\dfrac{2}{3}\,J^{-5/3}},
\end{eqnarray}
%
\begin{eqnarray}
J=\left( \det(\textbf{C})\right)^{\frac{1}{2}}\Rightarrow\dfrac{\partial\, J}{\partial\, \textbf{C}}=\underbrace{\dfrac{1}{2}\left( \det(\textbf{C})\right)^{-\frac{1}{2}}\dfrac{\partial\, \det(\textbf{C})}{\partial \,\textbf{C}}}_{\underbrace{\dfrac{\partial \det(\textbf{C})}{\partial \textbf{C}}=\det(\textbf{C})\,\textbf{C}^{-\text{T}}}_{\text{Classical identity.}}}=\underbrace{\dfrac{1}{2}\left( \det(\textbf{C})\right)^{\frac{1}{2}}\textbf{C}^{-\text{T}}}_{\textbf{C}=\textbf{C}^{\text{T}}\Rightarrow\textbf{C}^{-1}=\textbf{C}^{\text{-T}}}\Rightarrow\boxed{\dfrac{\partial \,J}{\partial\, \textbf{C}}=\dfrac{J}{2}\,\textbf{C}^{-1}},
\end{eqnarray}
%
\begin{eqnarray}
I_1=\tr(\textbf{C})=C_{11}+C_{22}+C_{33}\Rightarrow \dfrac{\partial\,I_1}{\partial\, C_{ij}}=\dfrac{\partial \,C_{kk}}{\partial\, C_{ij}}=\delta_{ij}\Rightarrow \boxed{\dfrac{\partial\, I_1}{\partial\,\textbf{C}}=\Id}.
\end{eqnarray}
%
Then the second Piola--Kirchhoff stress
$\textbf{S}$ is:
%
\begin{eqnarray}
\begin{array}{ccl}
\textbf{S}=2\dfrac{\partial \,\Psi}{\partial\,\textbf{C}}&=&
2\dfrac{\partial \left( \dfrac{\mu}{2}\left(J^{-2/3}I_1-3\right)
	+\dfrac{\kappa}{4}\left(J^2-1-2\ln J\right)\right) }{\partial\,\textbf{C}},\\[2ex]
&=&\mu\dfrac{\partial \,(J^{-2/3}I_1)}{\partial\,\textbf{C}}+\dfrac{\kappa}{2}\dfrac{\partial \,(J^{2}-2\,\ln J)}{\partial\,\textbf{C}},\\[2ex]
&=&\mu\left(\dfrac{\partial\,(J^{-2/3})}{\partial \,J}\dfrac{\partial \,J}{\partial\,\textbf{C}}\,I_1+J^{-2/3}\dfrac{\partial\,I_1}{\partial\,\textbf{C}}\right) +\dfrac{\kappa}{2}\dfrac{\partial\, (J^{2}-2\,\ln J)}{\partial\, J}\dfrac{\partial\, J}{\partial\,\textbf{C}},\\[2ex]
&=&\mu\left(\left(-\dfrac{2}{3}\,J^{-5/3}\right) \left(\dfrac{J}{2}\,\textbf{C}^{-1}\right) I_1+J^{-2/3}\textbf{I}\right) +\dfrac{\kappa}{2}\left(2\,J-\dfrac{2}{J} \right) \dfrac{J}{2}\,\textbf{C}^{-1},
	\end{array}
\end{eqnarray}
%
therefore
%
\begin{eqnarray}\label{eq_Sgeneral}
\boxed{\textbf{S}=\mu\,J^{-\frac{2}{3}}
\left(
\Id-\dfrac{I_1}{3}\textbf{C}^{-1}
\right)
+\dfrac{\kappa}{2}(J^2-1)\textbf{C}^{-1}=\mu\,J^{-\frac{2}{3}}
\,\Id
+\left( \dfrac{\kappa}{2}(J^2-1)-\mu\,J^{-\frac{2}{3}}\dfrac{I_1}{3}\right) \textbf{C}^{-1}}.
\end{eqnarray}
%
The first Piola--Kirchhoff stress is
%
\begin{equation}\label{eq_Pgeneral}
\boxed{\textbf{P}=\textbf{F}\textbf{S}
=\mu\, J^{-\frac{2}{3}}
\left(
\textbf{F}-\dfrac{I_1}{3}\textbf{F}^{-\text{T}}
\right)
+\dfrac{\kappa}{2}(J^2-1)\textbf{F}^{-\text{T}}=\mu\,J^{-\frac{2}{3}}
\,\textbf{F}
+\left( \dfrac{\kappa}{2}(J^2-1)-\mu\,J^{-\frac{2}{3}}\dfrac{I_1}{3}\right) \textbf{F}^{-\text{T}}}.
\end{equation}
%
The Kirchhoff stress is
%
\begin{equation}\label{eq_taugeneral}
	\boxed{\boldsymbol{\tau}=\textbf{P}\textbf{F}^{\text{T}}
	=\mu J^{-\frac{2}{3}}
	\left(
	\textbf{B}-\dfrac{I_1}{3}\Id
	\right)
	+\dfrac{\kappa}{2}(J^2-1)\Id=\mu\,J^{-\frac{2}{3}}
	\,\textbf{B}
	+\left( \dfrac{\kappa}{2}(J^2-1)-\mu\,J^{-\frac{2}{3}}\dfrac{I_1}{3}\right) \Id}.
\end{equation}
%
Finally,
%
\begin{equation}\label{eq_sigmaSource}
	\boxed{\boldsymbol{\sigma}=\dfrac1J\boldsymbol{\tau}
	=\mu J^{-\frac{5}{3}}
	\left(
	\textbf{B}-\dfrac{I_1}{3}\Id
	\right)
	+\dfrac{\kappa}{2}\left( \dfrac{J^2-1}{J}\right) \Id=\mu\,J^{-\frac{5}{3}}
	\,\textbf{B}
	+\left( \dfrac{\kappa}{2}\left( \dfrac{J^2-1}{J}\right)-\mu\,J^{-\frac{5}{3}}\dfrac{I_1}{3}\right) \Id}.
\end{equation}
%
Other expressions for $\boldsymbol{\tau}$ and $\boldsymbol{\sigma}$ are
%
\begin{equation}
	\boldsymbol{\tau}
=\mu\,\dev(\overline{\textbf{B}})
+\dfrac{\kappa}{2}\left( J^2-1\right) \Id
\quad\text{and}\quad
	\boldsymbol{\sigma}
	=\dfrac{\mu}{J}\dev(\overline{\textbf{B}})
	+\dfrac{\kappa}{2}\left( \dfrac{J^2-1}{J}\right) \Id,
\end{equation}
%
where
%
\begin{equation}
	\dev(\overline{\textbf{B}})=\overline{\textbf{B}}-\dfrac13\tr(\overline{\textbf{B}})\Id\quad\text{and}\quad \overline{\textbf{B}}=J^{-2/3}\textbf{B}.
\end{equation}

Defining
%
\begin{eqnarray}\label{eq_mc}
\boxed{m=\mu\,J^{-2/3}\quad\text{and}\quad c=\dfrac{\kappa}{2}\left(J^2-1\right)-m\,\dfrac{I_1}{3}},
\end{eqnarray}
%
then the compact form of stresses is:
%
\begin{eqnarray}\label{eq_stressesCompact}
\boxed{\textbf{S}=m\,\textbf{I}+c\,\textbf{C}^{-1},\quad\textbf{P}=m\,\textbf{F}+c\,\textbf{F}^{-\text{T}},\quad \boldsymbol{\tau}=m\,\textbf{B}+c\,\textbf{I}\quad\text{and}\quad \boldsymbol{\sigma}=\dfrac{m}{J}\,\textbf{B}+\dfrac{c}{J}\,\textbf{I}}.
\end{eqnarray}

\section{Manufactured displacement}

Suppose that the manufactured field has the form
%
\begin{equation}
\textbf{u}(\textbf{X},t)=\textbf{a}\,q(\textbf{X},t),
\label{eq_uCompact}
\end{equation}
%
where $\textbf{a}=[a_x,\,a_y,\,a_z]^{\text{T}}$ is constant, therefore $u_i=a_i\,q$. Define
%
\begin{equation}
\textbf{g}=\Grad_0 q,
\qquad
A=\textbf{a}\cdot\textbf{a},
\qquad
G=\textbf{g}\cdot\textbf{g}.
\end{equation}
%
Then
%
\begin{eqnarray}
\Grad_0\textbf{u}=\Grad_0(\textbf{a}\,q)\Rightarrow \dfrac{\partial u_i}{\partial X_j} = a_i\dfrac{\partial q}{\partial X_j}= a_i\,g_j\Rightarrow \boxed{\Grad_0\,\textbf{u}=\textbf{a}\,\textbf{g}^{\text{T}}=\textbf{a}\otimes\textbf{g}},
\end{eqnarray}
%
therefore
%
\begin{equation}
\boxed{\textbf{F}=\Id+\textbf{a}\,\textbf{g}^{\text{T}}=\Id+\textbf{a}\otimes\textbf{g}}.
\label{eq_FJcompact}
\end{equation}

For the determinant we can use the identity $\det(\textbf{A}+\textbf{u}\,\textbf{v}^\text{T})=\det(\textbf{A})(1+\textbf{v}^\text{T}\,\textbf{A}^{-1}\,\textbf{u})$, therefore:
%
\begin{eqnarray}
J=\det(\textbf{F})=\det(\Id+\textbf{a}\,\textbf{g}^\text{T})=\det(\Id)(1+\textbf{g}^\text{T}\,\Id\,\textbf{a})=1+\textbf{g}^\text{T}\,\textbf{a}\Rightarrow \boxed{J=1+\textbf{g}^\text{T}\,\textbf{a}=1+\textbf{a}\cdot\textbf{g}}.
\end{eqnarray}

The Sherman--Morrison identity is
%
\begin{eqnarray}
\left( \textbf{A}+\textbf{u}\,\textbf{v}^\text{T}\right) ^{-1}=\textbf{A}^{-1}-\dfrac{\textbf{A}^{-1}\,\textbf{u}\,\textbf{v}^\text{T}\,\textbf{A}^{-1}}{1+\textbf{v}^\text{T}\,\textbf{A}^{-1}\,\textbf{u}},
\end{eqnarray}
%
therefore:
%
\begin{eqnarray}
\textbf{F}^{-1}=\left( \Id+\textbf{a}\,\textbf{g}^{\text{T}}\right)^{-1}=\Id^{-1}-\dfrac{\Id^{-1}\,\textbf{a}\,\textbf{g}^\text{T}\,\Id^{-1}}{1+\textbf{g}^\text{T}\,\Id^{-1}\,\textbf{a}}=\Id-\dfrac{\textbf{a}\,\textbf{g}^\text{T}}{1+\textbf{g}^\text{T}\,\textbf{a}}=\Id-\dfrac{\textbf{a}\,\textbf{g}^\text{T}}{J}=\Id-\dfrac{\textbf{a}\otimes\textbf{g}}{J},
\end{eqnarray}
%
then
%
\begin{equation}
\boxed{
\textbf{F}^{-1}
=\Id-\dfrac{\textbf{a}\,\textbf{g}^\text{T}}{J}=\Id-\dfrac{\textbf{a}\otimes\textbf{g}}{J},
\qquad
\textbf{F}^{-\text{T}}
=\Id-\dfrac{\textbf{g}\,\textbf{a}^\text{T}}{J}=\Id-\dfrac{\textbf{g}\otimes\textbf{a}}{J}
}.
\label{eq_FinvCompact}
\end{equation}

The strain tensors are
%
\begin{eqnarray}
\textbf{C}=\textbf{F}^\text{T}\textbf{F}=\left( \Id+\textbf{g}\,\textbf{a}^{\text{T}}\right)\left( \Id+\textbf{a}\,\textbf{g}^{\text{T}}\right) =\Id+\textbf{g}\,\textbf{a}^{\text{T}}+\textbf{a}\,\textbf{g}^{\text{T}}+\textbf{g}\,\textbf{a}^{\text{T}}\,\textbf{a}\,\textbf{g}^{\text{T}}=\Id+\textbf{g}\,\textbf{a}^{\text{T}}+\textbf{a}\,\textbf{g}^{\text{T}}+\textbf{g}\,(\textbf{a}\cdot\textbf{a})\,\textbf{g}^{\text{T}}
\end{eqnarray}
%
\begin{eqnarray}
\textbf{B}=\textbf{F}\textbf{F}^\text{T}=\left( \Id+\textbf{a}\,\textbf{g}^{\text{T}}\right)\left( \Id+\textbf{g}\,\textbf{a}^{\text{T}}\right) =\Id+\textbf{a}\,\textbf{g}^{\text{T}}+\textbf{g}\,\textbf{a}^{\text{T}}+\textbf{a}\,\textbf{g}^{\text{T}}\,\textbf{g}\,\textbf{a}^{\text{T}}=\Id+\textbf{a}\,\textbf{g}^{\text{T}}+\textbf{g}\,\textbf{a}^{\text{T}}+\textbf{a}\,(\textbf{g}\cdot\textbf{g})\,\textbf{a}^{\text{T}}
\end{eqnarray}
%
therefore
%
\begin{equation}
\boxed{
\textbf{C}
=\Id+\textbf{g}\,\textbf{a}^{\text{T}}+\textbf{a}\,\textbf{g}^{\text{T}}+A\,\textbf{g}\,\textbf{g}^{\text{T}}=\Id+\textbf{g}\otimes\textbf{a}
+\textbf{a}\otimes\textbf{g}
+A\,\textbf{g}\otimes\textbf{g}
},
\label{eq_Ccompact}
\end{equation}
%
\begin{equation}
\boxed{
\textbf{B}
=\Id+\textbf{a}\,\textbf{g}^{\text{T}}+\textbf{g}\,\textbf{a}^{\text{T}}+G\,\textbf{a}\,\textbf{a}^{\text{T}}=\Id+\textbf{a}\otimes\textbf{g}
+\textbf{g}\otimes\textbf{a}
+G\,\textbf{a}\otimes\textbf{a}
},
\label{eq_Bcompact}
\end{equation}
%
and
%
\begin{eqnarray}
\begin{array}{rcl}
\textbf{C}^{-1}&=&\left(\textbf{F}^\text{T}\textbf{F}\right)^{-1} =\textbf{F}^{-1}\textbf{F}^{-\text{T}}=\left( \Id-\dfrac{\textbf{a}\,\textbf{g}^\text{T}}{J}\right) \left(\Id-\dfrac{\textbf{g}\,\textbf{a}^\text{T}}{J} \right) =\Id -\dfrac{\textbf{a}\,\textbf{g}^\text{T}}{J}-\dfrac{\textbf{g}\,\textbf{a}^\text{T}}{J}+\dfrac{\textbf{a}\,\textbf{g}^\text{T}}{J}\dfrac{\textbf{g}\,\textbf{a}^\text{T}}{J}\\
&=&\Id -\dfrac{\textbf{a}\,\textbf{g}^\text{T}}{J}-\dfrac{\textbf{g}\,\textbf{a}^\text{T}}{J}+\dfrac{1}{J^2}\textbf{a}\,\left( \textbf{g}\cdot\textbf{g}\right) \,\textbf{a}^\text{T}=\Id -\left( \dfrac{\textbf{a}\,\textbf{g}^\text{T}+\textbf{g}\,\textbf{a}^\text{T}}{J}\right) +\dfrac{\left( \textbf{g}\cdot\textbf{g}\right)}{J^2}\,\textbf{a}\, \textbf{a}^\text{T},
\end{array}
\end{eqnarray}
%
therefore
%
\begin{equation}
\boxed{
\textbf{C}^{-1}=
\Id -\left( \dfrac{\textbf{a}\,\textbf{g}^\text{T}+\textbf{g}\,\textbf{a}^\text{T}}{J}\right) +\dfrac{G}{J^2}\,\textbf{a}\,\textbf{a}^\text{T}=\Id
-\left( \dfrac{\textbf{a}\otimes\textbf{g}+\textbf{g}\otimes\textbf{a}}{J}\right) 
+\dfrac{G}{J^2}\,\textbf{a}\otimes\textbf{a}
}.
\label{eq_CinvCompact}
\end{equation}
%
In three dimensions, including the plane-strain embedding used later,
%
\begin{equation}
\boxed{
I_1=\tr(\textbf{C})=3+2\,\textbf{a}\cdot\textbf{g}+A\,G
}.
\label{eq_I1compact}
\end{equation}

Define also
\begin{equation}
\boxed{
d=m+c\quad\text{and}
\quad
e=\dfrac{c}{J}
},
\label{eq_de}
\end{equation}
%
and using Eqs.~\eqref{eq_stressesCompact},~\eqref{eq_FJcompact}~and~\eqref{eq_FinvCompact}, the first Piola stress becomes
%
\begin{equation}
\boxed{
\textbf{P}
=d\,\Id+m\,\textbf{a}\,\textbf{g}^\text{T}
-e\,\textbf{g}\,\textbf{a}^\text{T}=d\,\Id+m\,\textbf{a}\otimes\textbf{g}
-e\,\textbf{g}\otimes\textbf{a}
}.
\label{eq_PcompactRankOne}
\end{equation}

\section{Exact material divergence and manufactured body force}

For later use define
%
\begin{equation}
\boxed{\textbf{H}=\Grad_0\,\textbf{g},
\qquad
\textbf{h}=\textbf{H}\textbf{a}=\Grad_0\, J,
\qquad
\textbf{r}=\textbf{H}\textbf{g}},
\end{equation}
% 
and the following equalities
%
\begin{eqnarray}
\boxed{\Grad\cdot (\textbf{x}\,\textbf{y}^\text{T})=\left(\Grad\,\textbf{x}\right)\textbf{y}+\textbf{x}\left(\Grad\cdot\textbf{y} \right)},
\end{eqnarray}
%
\begin{eqnarray}
\boxed{\Grad(\textbf{x}\cdot\textbf{y})= \left( \Grad \textbf{x}\right)^\text{T}\textbf{y}+\left(\Grad\textbf{y} \right)^\text{T}\textbf{x}}.
\end{eqnarray}
%
\begin{eqnarray}
	\boxed{\Grad(\alpha\,\textbf{x})= \textbf{x}\left( \Grad \alpha\right)^\text{T}+\alpha\,\Grad\textbf{x}}.
\end{eqnarray}

Then the divergence of \textbf{P} is:
%
\begin{equation}
\begin{array}{rcl}
\Grad_0\cdot\textbf{P}
&=&\Grad_0\cdot\left( d\,\Id+m\,\textbf{a}\,\textbf{g}^\text{T}
-e\,\textbf{g}\,\textbf{a}^\text{T}\right),\\[1ex]
&=&\Grad_0\cdot\left( d\,\Id\right) +\Grad_0\cdot\left(m\,\textbf{a}\,\textbf{g}^\text{T}\right) 
-\Grad_0\cdot\left(e\,\textbf{g}\,\textbf{a}^\text{T}\right),
\\[1ex]
&=&\Grad_0\,d +\Grad_0\cdot\left(m\,\textbf{a}\,\textbf{g}^\text{T}\right) 
-\Grad_0\cdot\left(e\,\textbf{g}\,\textbf{a}^\text{T}\right),\\[1ex]
&=&\Grad_0\,d +\Grad_0\left(m\,\textbf{a}\right) \textbf{g} +m\,\textbf{a}\,\Grad_0\cdot \textbf{g}
-\Grad_0\left(e\,\textbf{g}\right) \textbf{a} -e\,\textbf{g}\,\underbrace{\Grad_0\cdot \textbf{a}}_{=0},\\[1ex]
&=&\Grad_0\,d +\textbf{a}\,\underbrace{\left(\Grad_0\, m\right)^\text{T} \textbf{g}}_{\textbf{g}\cdot \Grad_0\,m}+m\underbrace{\left( \Grad_0\, \textbf{a}\right)}_{=\textbf{0}}\textbf{g}  +m\,\textbf{a}\,\underbrace{\Grad_0\cdot \textbf{g}}_{=\Delta_0\, q}
-\textbf{g}\underbrace{\left(\Grad_0\, e\right)^\text{T} \textbf{a}}_{\textbf{a}\cdot \Grad_0\, e} -e\underbrace{\left( \Grad_0\, \textbf{g}\right) \textbf{a}}_{=\textbf{h}},
\end{array}
\end{equation}
%
Therefore:
%
\begin{equation}
\boxed{
\Grad_0\cdot\textbf{P}
=\Grad_0\, d
+\textbf{a}\left(\textbf{g}\cdot\Grad_0\, m+m\,\Delta_0\, q\right)
-\textbf{g}\left(\textbf{a}\cdot\Grad_0\, e\right)
-e\,\textbf{h}
}.
\label{eq_DivPCompact}
\end{equation}
%
Consequently,
%
\begin{equation}
\boxed{
\textbf{f}_0
=\rho_0\,\textbf{a}\,\ddot q
-\Grad_0\, d
-\textbf{a}\left(\textbf{g}\cdot\Grad_0\, m+m\,\Delta_0\, q\right)
+\textbf{g}\left(\textbf{a}\cdot\Grad_0\, e\right)
+e\,\textbf{h}
}.
\label{eq_fCompact}
\end{equation}
%
The rest of the expressions are:
%
\begin{eqnarray}\label{eq_gradm}
\Grad_0\, m=\Grad_0\left(\mu\,J^{-\frac{2}{3}}\right)=-\dfrac{2}{3}\mu\,J^{-\frac{5}{3}}\Grad_0\, J =-\dfrac{2\,\mu\,J^{-\frac{2}{3}}}{3\,J}\,\textbf{h}\Rightarrow\boxed{\Grad_0 m=-\dfrac{2\,m}{3\,J}\,\textbf{h}},
\end{eqnarray}
%
\begin{eqnarray}\label{eq_gradI1}
\begin{array}{rcl}
\Grad_0\, I_1&=&\Grad_0\left(3+2\,\textbf{a}\cdot\textbf{g}+A\,G\right)=2\,\underbrace{\left( \Grad_0\, \textbf{a}\right)^\text{T}}_{=\textbf{0}}\textbf{g}+2\,\underbrace{\left( \Grad_0\, \textbf{g}\right)^\text{T}\textbf{a}}_{\textbf{H}=\textbf{H}^\text{T}\Rightarrow\textbf{h}} +G\,\underbrace{\Grad_0 A}_{=\textbf{0}} + A\,\Grad_0 \left( \textbf{g}\cdot \textbf{g}\right),\\[1ex]
&=&2\,\textbf{h} + 2\,A\,\underbrace{\left( \Grad_0\, \textbf{g}\right)^\text{T}\textbf{g}}_{\textbf{H}=\textbf{H}^\text{T}\Rightarrow\textbf{r}}\Rightarrow \boxed{\Grad_0 I_1=2\textbf{h}+2A\textbf{r}}.
\end{array}
\end{eqnarray}
%
\begin{eqnarray}\label{eq_gradc}
\begin{array}{rcl}
\Grad_0\, c&=&\Grad_0\left(\dfrac{\kappa}{2}\,(J^2-1)-m\,\dfrac{I_1}{3}\right)=\kappa\,J\,\Grad_0\, J-\dfrac{I_1}{3}\Grad_0\, m - \dfrac{m}{3}\Grad_0\, I_1,\\[1ex] 
&&\Rightarrow\boxed{\Grad_0\, c=\kappa\,J\,\textbf{h}-\dfrac{I_1}{3}\Grad_0\, m - \dfrac{m}{3}\Grad_0\, I_1},
\end{array}
\end{eqnarray}
%
\begin{eqnarray}
\boxed{
\Grad_0\, d=\Grad_0\, m+\Grad_0\, c,
\qquad
\Grad_0\, e=\dfrac1J\Grad_0\, c-\dfrac{c}{J^2}\,\textbf{h}
}.
\label{eq_gradde}
\end{eqnarray}



\section{Two-dimensional problem: plane-strain use of the 3D constitutive law}

\subsection{Important interpretation of the 2D case}

The constitutive law in Eq.~\eqref{eq_sigmaSource} is a three-dimensional law:
it uses $J^{-2/3}$ and the three-dimensional deviatoric operator with the factor
$1/3$. Therefore, to use \emph{exactly the same material model} in a 2D problem,
we adopt plane strain,
%
\begin{equation}
\boxed{
u_z=0,
\qquad
\dfrac{\partial(\cdot)}{\partial z}=0,
\qquad
F_{zz}=1
}.
\end{equation}
%
The tensors are still three-dimensional, although only the in-plane momentum
equations are solved. In particular, $\sigma_{zz}$ is generally nonzero. Replacing
$J^{-2/3}$ by $J^{-1}$ or replacing $1/3$ by $1/2$ would define a different,
intrinsically two-dimensional constitutive model and would no longer reproduce
Eq.~\eqref{eq_sigmaSource}.

\subsection{Manufactured displacement and derivatives}

Let
%
\begin{equation}
q(x,y,t)=\sin(\omega_x\,x)\sin(\omega_y\,y)\sin(\omega_t\,t),
\end{equation}
%
with
%
\begin{equation}
\boxed{
\textbf{u}(x,y,t)=
\begin{bmatrix}a_x\,q\\a_y\,q\end{bmatrix}
}.                
\label{eq_u2D}    
\end{equation}
%
For the three-dimensional plane-strain embedding we use
%
\begin{equation}
\textbf{a}=\begin{bmatrix}a_x\\a_y\\0\end{bmatrix},
\qquad
\textbf{g}=\begin{bmatrix}g_x\\g_y\\0\end{bmatrix}.
\end{equation}
%
Introduce
%
\begin{eqnarray}
s_x=\sin(\omega_x\,x),\,\,\, c_x=\cos(\omega_x\,x),\,\,\,
s_y=\sin(\omega_y\,y),\,\,\, c_y=\cos(\omega_y\,y),\,\,\,s_t=\sin(\omega_t\,t),\,\,\, c_t=\cos(\omega_t\,t).
\end{eqnarray}
%
Then
%
\begin{equation}
\boxed{
q=s_x\,s_y\,s_t,
\quad
g_x=\omega_x\,c_x\,s_y\,s_t,
\quad
g_y=\omega_y\,s_x\,c_y\,s_t
}.
\end{equation}
%
The nonzero Hessian components are
%
\begin{equation}
\boxed{
H_{xx}=-\omega_x^2\,q,
\qquad
H_{yy}=-\omega_y^2\,q,
\qquad
H_{xy}=H_{yx}=\omega_x\,\omega_y\,c_x\,c_y\,s_t
}.
\end{equation}
%
Moreover,
%
\begin{equation}
\boxed{
\Delta\, q=-(\omega_x^2+\omega_y^2)\,q,
\qquad
\ddot q=-\omega_t^2\,q
}.
\label{eq_qDeriv2D}
\end{equation}

\subsection{Deformation and strain tensors in 2D plane strain}

The three-dimensional deformation gradient associated with the 2D field is
%
\begin{equation}
\boxed{
\textbf{F}=
\begin{bmatrix}
1+a_x\,g_x&a_x\,g_y&0\\
a_y\,g_x&1+a_y\,g_y&0\\
0&0&1
\end{bmatrix}
}.
\label{eq_F2D}
\end{equation}
%
Its determinant is
%
\begin{equation}
\boxed{J=1+a_x\,g_x+a_y\,g_y}.
\label{eq_J2D}
\end{equation}
%
The inverse transpose is
%
\begin{equation}
\boxed{
\textbf{F}^{-\text{T}}=
\begin{bmatrix}
\dfrac{1+a_y\,g_y}{J}&-\dfrac{a_y\,g_x}{J}&0\\[1.3ex]
-\dfrac{a_x\,g_y}{J}&\dfrac{1+a_x\,g_x}{J}&0\\[1.3ex]
0&0&1
\end{bmatrix}
}.
\label{eq_FinvT2D}
\end{equation}
%
Define
%
\begin{equation}
A=a_x^2+a_y^2,
\qquad
G=g_x^2+g_y^2.
\end{equation}
%
The right Cauchy--Green tensor is
%
\begin{equation}
\boxed{
\textbf{C}=
\begin{bmatrix}
1+2\,a_x\,g_x+A\,g_x^2&a_y\,g_x+a_x\,g_y+A\,g_x\,g_y&0\\
a_y\,g_x+a_x\,g_y+A\,g_x\,g_y&1+2\,a_y\,g_y+A\,g_y^2&0\\
0&0&1
\end{bmatrix}
}.
\label{eq_C2D}
\end{equation}
%
The Green--Lagrange strain is
%
\begin{equation}
\boxed{
\textbf{E}=\dfrac12
\begin{bmatrix}
2\,a_x\,g_x+A\,g_x^2&a_y\,g_x+a_x\,g_y+A\,g_x\,g_y&0\\
a_y\,g_x+a_x\,g_y+A\,g_x\,g_y&2\,a_y\,g_y+A\,g_y^2&0\\
0&0&0
\end{bmatrix}
}.
\label{eq_E2D}
\end{equation}
%
The left Cauchy--Green tensor is
%
\begin{equation}
\boxed{
\textbf{B}=
\begin{bmatrix}
1+2\,a_x\,g_x+G\,a_x^2&a_x\,g_y+a_y\,g_x+G\,a_x\,a_y&0\\
a_x\,g_y+a_y\,g_x+G\,a_x\,a_y&1+2\,a_y\,g_y+G\,a_y^2&0\\
0&0&1
\end{bmatrix}
}.
\label{eq_B2D}
\end{equation}
%
The isochoric left Cauchy--Green tensor is simply
%
\begin{equation}
\boxed{\overline{\textbf{B}}=J^{-2/3}\textbf{B}}.
\end{equation}
%
The first invariant is
%
\begin{equation}
\boxed{I_1=3+2\,(a_x\,g_x+a_y\,g_y)+A\,G}.
\label{eq_I12D}
\end{equation}
%
The inverse right Cauchy--Green tensor is
%
\begin{equation}
\boxed{
\textbf{C}^{-1}=
\begin{bmatrix}
D_{xx}&D_{xy}&0\\
D_{xy}&D_{yy}&0\\
0&0&1
\end{bmatrix}
},
\end{equation}
%
where
%
\begin{eqnarray}
D_{xx}&=&1-\dfrac{2\,a_x\,g_x}{J}+\dfrac{G\,a_x^2}{J^2},\\
D_{yy}&=&1-\dfrac{2\,a_y\,g_y}{J}+\dfrac{G\,a_y^2}{J^2},\\
D_{xy}&=&-\dfrac{a_x\,g_y+a_y\,g_x}{J}+\dfrac{G\,a_x\,a_y}{J^2}.
\end{eqnarray}

\subsection{Stress tensors in 2D}

As before, with $I_1$ given by Eq.~\eqref{eq_I12D},
%
\begin{equation}
m=\mu\,J^{-2/3},
\qquad
c=\dfrac{\kappa}{2}\,(J^2-1)-\dfrac{m\,I_1}{3},
\qquad
d=m+c,
\qquad
e=\dfrac{c}{J}.
\end{equation}
%
The second Piola--Kirchhoff stress is
%
\begin{equation}
\boxed{
\textbf{S}=\begin{bmatrix}
m+c\,D_{xx}&c\,D_{xy}&0\\
c\,D_{xy}&m+c\,D_{yy}&0\\
0&0&d
\end{bmatrix}
}.
\label{eq_S2D}
\end{equation}
%
The first Piola--Kirchhoff stress is
%
\begin{equation}
\boxed{
\textbf{P}=\begin{bmatrix}
d+(m-e)\,a_x\,g_x&m\, a_x\,g_y-e\, g_x\,a_y&0\\
m\, a_y\,g_x-e\, g_y\,a_x&d+(m-e)\,a_y\,g_y&0\\
0&0&d
\end{bmatrix}
}.
\label{eq_P2D}
\end{equation}
%
The Kirchhoff stress is
%
\begin{equation}
\boxed{
\boldsymbol{\tau}=\begin{bmatrix}
m\,B_{xx}+c&m\,B_{xy}&0\\
m\,B_{xy}&m\,B_{yy}+c&0\\
0&0&d
\end{bmatrix}
},
\label{eq_tau2D}
\end{equation}
%
where $B_{xx}$, $B_{xy}$ and $B_{yy}$ are the in-plane components of $\textbf{B}$ in Eq.~\eqref{eq_B2D}. The Cauchy stress is
%
\begin{equation}
\boxed{
\boldsymbol{\sigma}=\dfrac1J
\begin{bmatrix}
m\,B_{xx}+c&m\,B_{xy}&0\\
m\,B_{xy}&m\,B_{yy}+c&0\\
0&0&d
\end{bmatrix}
}.
\label{eq_sigma2D}
\end{equation}
%
The out-of-plane components $S_{zz}=P_{zz}=\tau_{zz}=d$ and $\sigma_{zz}=d/J$ are generally nonzero, but they do not depend on $z$ and therefore do not enter the in-plane momentum equations.


\subsection{Exact body force in 2D}

The vectors $\textbf{h}=\textbf{H}\textbf{a}$ and $\textbf{r}=\textbf{H}\textbf{g}$ have in-plane components
%
\begin{align}
h_x&=H_{xx}\,a_x+H_{xy}\,a_y,
&h_y&=H_{xy}\,a_x+H_{yy}\,a_y,
\\
r_x&=H_{xx}\,g_x+H_{xy}\,g_y,
&r_y&=H_{xy}\,g_x+H_{yy}\,g_y.
\end{align}
%
Then, for $i=x,y$, it is possible to obtain $(\bullet)_{,i}=\dfrac{\partial\,(\bullet)}{\partial\,i}$ as
%
\begin{eqnarray}
(I_1)_{,i}&=&2\,h_i+2\,A\,r_i,
\\
m_{,i}&=&-\dfrac{2\,m}{3\,J}\,h_i,
\\
c_{,i}&=&\kappa \,J\,h_i-\dfrac{I_1}{3}\,m_{,i}
-\dfrac{m}{3}\,(I_1)_{,i},
\\
d_{,i}&=&m_{,i}+c_{,i},
\\
e_{,i}&=&\dfrac{c_{,i}}{J}-\dfrac{c}{J^2}\,h_i.
\end{eqnarray}
%
Define the two scalars
%
\begin{equation}
\boxed{
R_{2D}=m\,\Delta q+g_x\,m_{,x}+g_y\,m_{,y},
\qquad
T_{2D}=a_x\,e_{,x}+a_y\,e_{,y}
}.
\end{equation}
%
The exact body force per unit reference volume is then $\textbf{f}_0=[f_{0x},\,f_{0y},\,f_{0z}]^\text{T}$:
%
\begin{equation}
\boxed{
\begin{aligned}
f_{0x}&=-\rho_0\,\omega_t^2\,a_x\,q-d_{,x}-a_x\,R_{2D}+g_x\,T_{2D}+e\,h_x,
\\
f_{0y}&=-\rho_0\,\omega_t^2\,a_y\,q-d_{,y}-a_y\,R_{2D}+g_y\,T_{2D}+e\,h_y,
\\
f_{0z}&=0
\end{aligned}
}.
\label{eq_f2D}
\end{equation}

\subsection{Dirichlet boundary conditions in 2D}

For
%
\begin{equation}
\Omega=[0,L_x]\times[0,L_y],
\end{equation}
%
the exact displacement is prescribed on every side:
%
\begin{align}
\textbf{u}(0,y,t)&=\textbf{0},\\
\textbf{u}(L_x,y,t)&=
\begin{bmatrix}a_x\\a_y\end{bmatrix}
\sin(\omega_x\,L_x)\sin(\omega_y\,y)\sin(\omega_t\,t),\\
\textbf{u}(x,0,t)&=\textbf{0},\\
\textbf{u}(x,L_y,t)&=
\begin{bmatrix}a_x\\a_y\end{bmatrix}
\sin(\omega_x\,x)\sin(\omega_y\,L_y)\sin(\omega_t\,t).
\end{align}
%
If
%
\begin{equation}
\omega_x=\dfrac{n_x\,\pi}{L_x},
\qquad
\omega_y=\dfrac{n_y\,\pi}{L_y},
\qquad n_x,n_y\in\mathbb{N},
\end{equation}
%
all boundary conditions are homogeneous Dirichlet conditions.

\subsection{Initial conditions in 2D}

Because $\sin(\omega_t t)=0$ at $t=0$,
%
\begin{equation}
\boxed{\textbf{u}(x,y,0)=\textbf{0}}.
\label{eq_u0_2D}
\end{equation}
%
The initial velocity is nonzero:
%
\begin{equation}
\boxed{
\dot{\textbf{u}}(x,y,0)
=\omega_t
\begin{bmatrix}a_x\\a_y\end{bmatrix}
\sin(\omega_x\,x)\sin(\omega_y\,y)
}.
\label{eq_v0_2D}
\end{equation}
%
Also $\ddot{\textbf{u}}(x,y,0)=\textbf{0}$. At the initial instant
%
\begin{equation}
\textbf{F}=\Id,
\qquad J=1,
\qquad \textbf{C}=\textbf{B}=\Id,
\end{equation}
%
and the constitutive law is stress free:
%
\begin{equation}
\boxed{
\textbf{S}=\textbf{P}=\boldsymbol{\tau}=\boldsymbol{\sigma}=\textbf{0}
\qquad (t=0)
}.
\end{equation}

\section{Three-dimensional manufactured solution}

\subsection{Displacement field and derivatives}

Let
%
\begin{equation}
q(x,y,z,t)=
\sin(\omega_x\,x)\sin(\omega_y\,y)\sin(\omega_z\,z)\sin(\omega_t\,t),
\end{equation}
%
and
%
\begin{equation}
\boxed{
\textbf{u}(x,y,z,t)=
\begin{bmatrix}a_x\,q\\a_y\,q\\a_z\,q\end{bmatrix}
}.
\label{eq_u3D}
\end{equation}
%
Define
%
\begin{equation}
s_z=\sin(\omega_z\,z),\qquad c_z=\cos(\omega_z\,z),
\end{equation}
%
with $s_x,c_x,s_y,c_y,s_t,c_t$ defined as before. Then
%
\begin{equation}
\boxed{
\textbf{g}=
\begin{bmatrix}
\omega_x\,c_x\,s_y\,s_z\,s_t\\
\omega_y\,s_x\,c_y\,s_z\,s_t\\
\omega_z\,s_x\,s_y\,c_z\,s_t
\end{bmatrix}
}.
\label{eq_g3D}
\end{equation}
%
The Hessian is
%
\begin{equation}
\boxed{
\textbf{H}=
\begin{bmatrix}
-\omega_x^2\,q&H_{xy}&H_{xz}\\
H_{xy}&-\omega_y^2\,q&H_{yz}\\
H_{xz}&H_{yz}&-\omega_z^2\,q
\end{bmatrix}
},
\label{eq_H3D}
\end{equation}
%
where
%
\begin{align}
H_{xy}&=\omega_x\,\omega_y\,c_x\,c_y\,s_z\,s_t,\\
H_{xz}&=\omega_x\,\omega_z\,c_x\,s_y\,c_z\,s_t,\\
H_{yz}&=\omega_y\,\omega_z\,s_x\,c_y\,c_z\,s_t.
\end{align}
%
Moreover,
%
\begin{equation}
\boxed{
\Delta\, q=-(\omega_x^2+\omega_y^2+\omega_z^2)\,q,
\qquad
\ddot q=-\omega_t^2\,q
}.
\label{eq_qDeriv3D}
\end{equation}

\subsection{Deformation and strain tensors in 3D}

The deformation gradient is
%
\begin{equation}
\boxed{
\textbf{F}=\begin{bmatrix}
1+a_x\,g_x&a_x\,g_y&a_x\,g_z\\
a_y\,g_x&1+a_y\,g_y&a_y\,g_z\\
a_z\,g_x&a_z\,g_y&1+a_z\,g_z
\end{bmatrix}
}.
\label{eq_F3D}
\end{equation}
%
Its Jacobian is
%
\begin{equation}
\boxed{J=1+a_x\,g_x+a_y\,g_y+a_z\,g_z}.
\label{eq_J3D}
\end{equation}
%
The inverse transpose is
%
\begin{equation}
\boxed{
\textbf{F}^{-\text{T}}=\Id-
\dfrac1J
\begin{bmatrix}
g_x\,a_x&g_x\,a_y&g_x\,a_z\\
g_y\,a_x&g_y\,a_y&g_y\,a_z\\
g_z\,a_x&g_z\,a_y&g_z\,a_z
\end{bmatrix}
}.
\end{equation}
%
Define
%
\begin{equation}
A=a_x^2+a_y^2+a_z^2,
\qquad
G=g_x^2+g_y^2+g_z^2.
\end{equation}
%
The right Cauchy--Green tensor is
%
\begin{equation}
\boxed{
\textbf{C}=
\begin{bmatrix}
C_{xx}&C_{xy}&C_{xz}\\
C_{xy}&C_{yy}&C_{yz}\\
C_{xz}&C_{yz}&C_{zz}
\end{bmatrix}
},
\end{equation}
%
where
%
\begin{align}
C_{xx}&=1+2\,a_x\,g_x+A\,g_x^2,
&C_{xy}&=a_y\,g_x+a_x\,g_y+A\,g_x\,g_y,
\\
C_{yy}&=1+2\,a_y\,g_y+A\,g_y^2,
&C_{xz}&=a_z\,g_x+a_x\,g_z+A\,g_x\,g_z,
\\
C_{zz}&=1+2\,a_z\,g_z+A\,g_z^2,
&C_{yz}&=a_z\,g_y+a_y\,g_z+A\,g_y\,g_z.
\end{align}
%
The Green--Lagrange strain is
%
\begin{equation}
\boxed{\textbf{E}=\dfrac12(\textbf{C}-\Id)}.
\end{equation}
%
The left Cauchy--Green tensor is
%
\begin{equation}
\boxed{
\textbf{B}=
\begin{bmatrix}
B_{xx}&B_{xy}&B_{xz}\\
B_{xy}&B_{yy}&B_{yz}\\
B_{xz}&B_{yz}&B_{zz}
\end{bmatrix}
},
\end{equation}
%
where
%
\begin{align}
B_{xx}&=1+2\,a_x\,g_x+G\,a_x^2,
&B_{xy}&=a_x\,g_y+a_y\,g_x+G\,a_x\,a_y,
\\
B_{yy}&=1+2\,a_y\,g_y+G\,a_y^2,
&B_{xz}&=a_x\,g_z+a_z\,g_x+G\,a_x\,a_z,
\\
B_{zz}&=1+2\,a_z\,g_z+G\,a_z^2,
&B_{yz}&=a_y\,g_z+a_z\,g_y+G\,a_y\,a_z.
\end{align}
%
The isochoric tensor is
%
\begin{equation}
\boxed{\overline{\textbf{B}}=J^{-2/3}\textbf{B}},
\end{equation}
%
and the first invariant is
%
\begin{equation}
\boxed{I_1=3+2\,(a_x\,g_x+a_y\,g_y+a_z\,g_z)+A\,G}.
\end{equation}
%
The inverse right Cauchy--Green tensor is
%
\begin{equation}
\boxed{
\textbf{C}^{-1}=
\begin{bmatrix}
D_{xx}&D_{xy}&D_{xz}\\
D_{xy}&D_{yy}&D_{yz}\\
D_{xz}&D_{yz}&D_{zz}
\end{bmatrix}
},
\end{equation}
%
with
%
\begin{align}
D_{xx}&=1-\dfrac{2\,a_x\,g_x}{J}+\dfrac{G\,a_x^2}{J^2},
&D_{xy}&=-\dfrac{a_x\,g_y+a_y\,g_x}{J}+\dfrac{G\,a_x\,a_y}{J^2},
\\
D_{yy}&=1-\dfrac{2\,a_y\,g_y}{J}+\dfrac{G\,a_y^2}{J^2},
&D_{xz}&=-\dfrac{a_x\,g_z+a_z\,g_x}{J}+\dfrac{G\,a_x\,a_z}{J^2},
\\
D_{zz}&=1-\dfrac{2\,a_z\,g_z}{J}+\dfrac{G\,a_z^2}{J^2},
&D_{yz}&=-\dfrac{a_y\,g_z+a_z\,g_y}{J}+\dfrac{G\,a_y\,a_z}{J^2}.
\end{align}

\subsection{Stress tensors in 3D}

As before,
%
\begin{equation}
m=\mu\,J^{-2/3},
\qquad
c=\dfrac{\kappa}{2}\,(J^2-1)-\dfrac{m\,I_1}{3},
\qquad
d=m+c,
\qquad
e=\dfrac{c}{J}.
\end{equation}
%
The second Piola--Kirchhoff stress is
%
\begin{equation}
\boxed{
\textbf{S}=\begin{bmatrix}
m+c\,D_{xx}&c\,D_{xy}&c\,D_{xz}\\
c\,D_{xy}&m+c\,D_{yy}&c\,D_{yz}\\
c\,D_{xz}&c\,D_{yz}&m+c\,D_{zz}
\end{bmatrix}
}.
\label{eq_S3D}
\end{equation}
%
The first Piola--Kirchhoff stress is
%
\begin{equation}
\boxed{
\textbf{P}=\begin{bmatrix}
d+(m-e)\,a_x\,g_x&m\, a_x\,g_y-e\, g_x\,a_y&m\, a_x\,g_z-e\, g_x\,a_z\\
m\, a_y\,g_x-e\, g_y\,a_x&d+(m-e)\,a_y\,g_y&m\, a_y\,g_z-e\, g_y\,a_z\\
m\, a_z\,g_x-e\, g_z\,a_x&m\, a_z\,g_y-e\, g_z\,a_y&d+(m-e)\,a_z\,g_z
\end{bmatrix}
}.
\label{eq_P3D}
\end{equation}
%
The Kirchhoff stress is
%
\begin{equation}
\boxed{
\boldsymbol{\tau}=\begin{bmatrix}
m\,B_{xx}+c&m\,B_{xy}&m\,B_{xz}\\
m\,B_{xy}&m\,B_{yy}+c&m\,B_{yz}\\
m\,B_{xz}&m\,B_{yz}&m\,B_{zz}+c
\end{bmatrix}
}.
\label{eq_tau3D}
\end{equation}
%
The Cauchy stress is
%
\begin{equation}
\boxed{
\boldsymbol{\sigma}=\dfrac1J
\begin{bmatrix}
m\,B_{xx}+c&m\,B_{xy}&m\,B_{xz}\\
m\,B_{xy}&m\,B_{yy}+c&m\,B_{yz}\\
m\,B_{xz}&m\,B_{yz}&m\,B_{zz}+c
\end{bmatrix}
}.
\label{eq_sigma3D}
\end{equation}

\subsection{Exact body force in 3D}

The vector $\textbf{h}=\textbf{H}\textbf{a}$ has components
%
\begin{align}
h_x&=-\omega_x^2\,q\,a_x+H_{xy}\,a_y+H_{xz}\,a_z,\\
h_y&=H_{xy}\,a_x-\omega_y^2\,q\,a_y+H_{yz}\,a_z,\\
h_z&=H_{xz}\,a_x+H_{yz}\,a_y-\omega_z^2\,q\,a_z.
\end{align}
%
The vector $\textbf{r}=\textbf{H}\textbf{g}$ has components
%
\begin{align}
r_x&=-\omega_x^2\,q\,g_x+H_{xy}\,g_y+H_{xz}\,g_z,\\
r_y&=H_{xy}\,g_x-\omega_y^2\,q\,g_y+H_{yz}\,g_z,\\
r_z&=H_{xz}\,g_x+H_{yz}\,g_y-\omega_z^2\,q\,g_z.
\end{align}
%
For $i=x,y,z$, it is possible to obtain $(\bullet)_{,i}=\dfrac{\partial\,(\bullet)}{\partial\,i}$ as
%
\begin{align}
(I_1)_{,i}&=2\,h_i+2\,A\,r_i,
\\
m_{,i}&=-\dfrac{2\,m}{3\,J}\,h_i,
\\
c_{,i}&=\kappa\, J\,h_i-\dfrac{I_1}{3}\,m_{,i}
-\dfrac{m}{3}\,(I_1)_{,i},
\\
d_{,i}&=m_{,i}+c_{,i},
\\
e_{,i}&=\dfrac{c_{,i}}{J}-\dfrac{c}{J^2}\,h_i.
\end{align}
%
Define
%
\begin{equation}
\boxed{
R_{3D}=m\,\Delta q+g_x\,m_{,x}+g_y\,m_{,y}+g_z\,m_{,z}
},
\end{equation}
%
and
%
\begin{equation}
\boxed{
T_{3D}=a_x\,e_{,x}+a_y\,e_{,y}+a_z\,e_{,z}
}.
\end{equation}
%
Then the exact body-force vector is
%
\begin{equation}
\boxed{
\begin{aligned}
f_{0x}&=-\rho_0\,\omega_t^2\,a_x\,q-d_{,x}-a_x\,R_{3D}+g_x\,T_{3D}+e\,h_x,
\\
f_{0y}&=-\rho_0\,\omega_t^2\,a_y\,q-d_{,y}-a_y\,R_{3D}+g_y\,T_{3D}+e\,h_y,
\\
f_{0z}&=-\rho_0\,\omega_t^2\,a_z\,q-d_{,z}-a_z\,R_{3D}+g_z\,T_{3D}+e\,h_z
\end{aligned}
}.
\label{eq_f3D}
\end{equation}


\subsection{Dirichlet boundary conditions in 3D}

For
%
\begin{equation}
\Omega=[0,L_x]\times[0,L_y]\times[0,L_z],
\end{equation}
%
$\textbf{u}=\textbf{u}^{\text{ex}}$ is imposed on all six faces:
%
\begin{align}
\textbf{u}(0,y,z,t)&=\textbf{0},\\
\textbf{u}(L_x,y,z,t)&=\textbf{a}\sin(\omega_x\,L_x)\sin(\omega_y\,y)
\sin(\omega_z\,z)\sin(\omega_t\,t),\\
\textbf{u}(x,0,z,t)&=\textbf{0},\\
\textbf{u}(x,L_y,z,t)&=\textbf{a}\sin(\omega_x\,x)\sin(\omega_y\,L_y)
\sin(\omega_z\,z)\sin(\omega_t\,t),\\
\textbf{u}(x,y,0,t)&=\textbf{0},\\
\textbf{u}(x,y,L_z,t)&=\textbf{a}\sin(\omega_x\,x)\sin(\omega_y\,y)
\sin(\omega_z\,L_z)\sin(\omega_t\,t).
\end{align}
%
If
%
\begin{equation}
\omega_i=\dfrac{n_i\,\pi}{L_i},
\qquad n_i\in\mathbb{N},
\qquad i=x,y,z,
\end{equation}
%
all six conditions are homogeneous Dirichlet conditions.

\subsection{Initial conditions in 3D}

The exact initial displacement is
%
\begin{equation}
\boxed{\textbf{u}(x,y,z,0)=\textbf{0}}.
\label{eq_u0_3D}
\end{equation}
%
The initial velocity is
%
\begin{equation}
\boxed{
\dot{\textbf{u}}(x,y,z,0)
=\omega_t\,\textbf{a}\,
\sin(\omega_x\,x)\sin(\omega_y\,y)\sin(\omega_z\,z)
}.
\label{eq_v0_3D}
\end{equation}
%
Also $\ddot{\textbf{u}}(x,y,z,0)=\textbf{0}$. Since $\textbf{g}=\textbf{0}$ at $t=0$,
%
\begin{equation}
\textbf{F}=\Id,
\qquad
J=1,
\qquad
\textbf{C}=\textbf{B}=\Id,
\end{equation}
%
and therefore
%
\begin{equation}
\boxed{
\textbf{S}=\textbf{P}=\boldsymbol{\tau}=\boldsymbol{\sigma}=\textbf{0}
\qquad (t=0)
}.
\end{equation}

\section{Implementation summary}

For each point $\textbf{X}$ and time $t$, a direct implementation of the MMS can follow this sequence:
\begin{enumerate}
%
\item Evaluate $q$, $\textbf{g}=\Grad_0 q$, $\textbf{H}=\Grad_0\textbf{g}$, $\Delta_0 q$, and $\ddot q$: Eqs.~\eqref{eq_u2D}--\eqref{eq_qDeriv2D} in 2D and Eqs.~\eqref{eq_u3D}--\eqref{eq_qDeriv3D} in 3D. In 2D, use the plane-strain embedding $a_z=g_z=0$ and keep all tensors $3\times3$.
%
\item Evaluate $\textbf{F}=\Id+\textbf{a}\otimes\textbf{g}$ and $J=1+\textbf{a}\cdot\textbf{g}$ from Eq.~\eqref{eq_FJcompact} and Eqs.~\eqref{eq_J2D} or~\eqref{eq_J3D}. Check that $J>0$.
%
\item Evaluate $A=\textbf{a}\cdot\textbf{a}$, $G=\textbf{g}\cdot\textbf{g}$, and $I_1=3+2\,\textbf{a}\cdot\textbf{g}+A\,G$ from Eq.~\eqref{eq_I1compact}. The constant is $3$ also in 2D (Eq.~\eqref{eq_I12D}), because $C_{zz}=1$.
%
\item Evaluate the scalars of Eqs.~\eqref{eq_mc} and~\eqref{eq_de},
%
$$
m=\mu\,J^{-2/3},
\qquad
c=\dfrac{\kappa}{2}(J^2-1)-\dfrac{m\,I_1}{3},
\qquad
d=m+c,
\qquad
e=\dfrac{c}{J}.
$$
\item Construct the required stress measure with Eq.~\eqref{eq_stressesCompact},
%
$$
\textbf{S}=m\,\Id+c\,\textbf{C}^{-1},\quad
\textbf{P}=m\,\textbf{F}+c\,\textbf{F}^{-\text{T}},\quad
\boldsymbol{\tau}=m\,\textbf{B}+c\,\Id,\quad
\boldsymbol{\sigma}=J^{-1}\boldsymbol{\tau},
$$
%
where $\textbf{F}^{-\text{T}}$, $\textbf{B}$ and $\textbf{C}^{-1}$ are given by Eqs.~\eqref{eq_FinvCompact}, \eqref{eq_Bcompact} and~\eqref{eq_CinvCompact}. Equivalently, $\textbf{P}=d\,\Id+m\,\textbf{a}\otimes\textbf{g}-e\,\textbf{g}\otimes\textbf{a}$ (Eq.~\eqref{eq_PcompactRankOne}). In components, see Eqs.~\eqref{eq_S2D}--\eqref{eq_sigma2D} in 2D and Eqs.~\eqref{eq_S3D}--\eqref{eq_sigma3D} in 3D.
%
\item Evaluate $\textbf{h}=\textbf{H}\textbf{a}=\Grad_0 J$, $\textbf{r}=\textbf{H}\textbf{g}$, and the scalar gradients $\Grad_0 m$, $\Grad_0 I_1$, $\Grad_0 c$, $\Grad_0 d$ and $\Grad_0 e$ from Eqs.~\eqref{eq_gradm}--\eqref{eq_gradde}.
%
\item Evaluate $\Grad_0\cdot\textbf{P}$ from Eq.~\eqref{eq_DivPCompact} and then $\textbf{f}_0$ from Eq.~\eqref{eq_fCompact}. In components, see Eq.~\eqref{eq_f2D} in 2D (where $f_{0z}=0$) and Eq.~\eqref{eq_f3D} in 3D.
%
\item Impose $\textbf{u}=\textbf{u}^{\text{ex}}$ on the complete boundary, Eq.~\eqref{eq_allDirichlet}, and use the exact initial displacement and velocity, Eqs.~\eqref{eq_u0_2D}--\eqref{eq_v0_2D} in 2D and Eqs.~\eqref{eq_u0_3D}--\eqref{eq_v0_3D} in 3D.
%
\end{enumerate}

The body force in this report is per unit \emph{reference} volume. If a solver instead uses the spatial equation
%
\begin{equation}
\rho\,\ddot{\textbf{u}}=\nabla_{\textbf{x}}\cdot\boldsymbol{\sigma}+\textbf{f},
\end{equation}
%
then $\rho=\rho_0/J$ and the corresponding body force per unit current volume is
%
\begin{equation}
\boxed{\textbf{f}=\dfrac{1}{J}\,\textbf{f}_0}.
\end{equation}

\end{document}
